\documentclass[conference]{article}
\usepackage[utf8]{inputenc}
\usepackage[T1]{fontenc}
\usepackage{amsmath,amssymb}
\usepackage{booktabs}
\usepackage{multirow}
\usepackage{graphicx}
\usepackage{xcolor}
\usepackage[margin=1in]{geometry}
\usepackage{hyperref}
\usepackage{tabularx}
\usepackage{caption}

\newcommand{\ours}{Meta-Brain}
\newcommand{\liberoplus}{LIBERO-Plus}

\title{Meta-Brain: Perturbation-Diagnostic Routing of LoRA Experts\\for Parameter-Efficient VLA Robustness}

\author{
  Anonymous Authors\\
  Submitted to CoRL 2026
}

\date{}

\begin{document}
\maketitle

% ============================================================
\begin{abstract}
Vision-Language-Action models (VLAs) achieve over 90\% success on standard manipulation benchmarks but collapse under environmental perturbations---frozen pi0 drops to 53.6\% on \liberoplus{}.
Existing robustness methods either apply perturbation-agnostic uniform adaptation or require prohibitive parameter budgets (e.g., VLA-GSE: 114M trainable parameters).
We observe that VLA internal representations already encode perturbation type with near-perfect accuracy: a linear probe on PaliGemma layer~6 activations achieves 3-class macro-F1$=$0.993 in a single forward pass.
Building on this finding, we propose \ours{}, a modular framework that \emph{diagnoses} the perturbation type from frozen activations via a lightweight Perturbation-Aware Router (PAR, 48K parameters), then \emph{routes} to one of three type-specialized LoRA experts (Visual, Spatial, Physical) for targeted correction.
On \liberoplus{}, \ours{} achieves 78.4\% with 6.1M trainable parameters, improving 24.8\% over the frozen baseline and 13.6\% over a parameter-matched universal LoRA ($p<0.001$).
With an optional Correction Memory module, performance reaches 80.8\%, closing 90\% of the gap to VLA-GSE (81.2\%, 114M params) at 19$\times$ fewer trainable parameters.
Cross-backbone validation on OpenVLA confirms the generalizability of perturbation-diagnostic routing (gain over universal LoRA: +11.8\%).
\end{abstract}

% ============================================================
\section{Introduction}

Vision-Language-Action models (VLAs) have demonstrated strong performance on standard manipulation benchmarks, yet they remain brittle when deployed in environments that differ from their training distribution.
On \liberoplus{}---a robustness benchmark that introduces six types of environmental perturbations---frozen pi0 drops from $>$90\% to 53.6\% success rate.

Different perturbation types disrupt different processing stages within the VLA.
Visual perturbations (camera angle, lighting, background clutter, sensor noise) corrupt early visual representations.
Spatial perturbations (object layout changes) invalidate mid-level spatial reasoning.
Physical perturbations (robot dynamics changes) break late-stage action planning.
Yet existing robustness methods treat all perturbations uniformly: VLA-GSE~\cite{vla-gse} routes by per-token capacity needs without perturbation awareness, requiring 114M trainable parameters; FTM~\cite{ftm} effectively corrects visual perturbations but cannot address spatial or physical failures; RobustVLA~\cite{robustvla} applies uniform adversarial training across all perturbation types.

We make three observations that motivate a new approach.
First, VLA internal activations already encode perturbation type: a linear probe on a single layer of frozen pi0 (PaliGemma layer~6) classifies perturbation type with F1$=$0.993 (\textbf{Claim~1}: perturbation diagnosability).
Second, type-specialized LoRA experts dramatically outperform a universal LoRA of equivalent total capacity: routed correction yields +13.6\% over uniform correction ($p<0.001$), indicating that the right fix for the right failure is fundamentally more efficient than a one-size-fits-all approach (\textbf{Claim~2}: targeted $>$ uniform correction).
Third, by leveraging this diagnostic signal for lightweight routing, \ours{} achieves competitive robustness with VLA-GSE at 19$\times$ fewer parameters (\textbf{Claim~3}: parameter-efficient robustness).

Our contributions are:
\begin{enumerate}
  \item An empirical finding that VLA internal representations encode perturbation type at single-layer, single-step granularity, enabling real-time diagnosis without external modules.
  \item \ours{}: a modular framework comprising a Perturbation-Aware Router (PAR) and an Adaptive LoRA Bank (ALB) that routes corrective experts by diagnosed perturbation type, with an optional Correction Memory (CM) module.
  \item A systematic design-dimension analysis comparing perturbation-diagnostic routing against capacity-based routing and universal adaptation, establishing perturbation-awareness as an orthogonal and composable design axis.
\end{enumerate}

% ============================================================
\section{Related Work}

\paragraph{VLA Robustness.}
\liberoplus{}~\cite{libero-plus} introduced a systematic perturbation benchmark for VLAs covering visual, spatial, and physical perturbation types.
SAFE~\cite{safe} demonstrated that VLA internal features can detect distribution shifts but did not close the loop with correction.
VLA-GSE~\cite{vla-gse} achieves 81.2\% on \liberoplus{} via capacity-based MoE routing with 114M trainable parameters on Qwen3-VL-4B.
RobustVLA~\cite{robustvla} applies adversarial training for +12.6\% uniform robustness gain.
None of these methods combine perturbation diagnosis with type-specific correction.

\paragraph{Parameter-Efficient Adaptation.}
LoRA~\cite{lora} and its variants provide parameter-efficient fine-tuning by injecting low-rank updates into frozen model layers.
Recent work has explored task-specific routing of LoRA modules: MoIRA and CORAL route by task identity, while LoRA-SP routes by skill category.
FTM/FLA~\cite{ftm} modulates visual features with extremely few parameters (4K) but is limited to visual perturbations.
Our work introduces a new routing dimension---perturbation type---that is orthogonal to task or capacity routing.

% ============================================================
\section{Method}

\ours{} operates in two phases on top of a frozen VLA backbone (pi0 with PaliGemma vision encoder).

\subsection{Phase 0: Perturbation Diagnosis}

We extract activations from a single intermediate layer of the frozen vision encoder (PaliGemma layer~6) during the standard forward pass.
A linear probe trained on these activations classifies the current observation into one of three perturbation categories: \textbf{Visual} (camera, lighting, background, noise), \textbf{Spatial} (layout), or \textbf{Physical} (robot dynamics).

This classification achieves 3-class macro-F1$=$0.993 on held-out data, confirming that the perturbation-type signal is already present in the frozen activations and does not require any additional forward passes or external detectors.
The mechanistic basis aligns with the VLA's layered information processing: visual perturbations primarily disrupt early layers, spatial perturbations affect mid-level representations, and physical perturbations impact late-stage action planning.

\subsection{Phase 1: Meta-Brain Architecture}

\paragraph{Perturbation-Aware Router (PAR).}
The PAR is a 2-layer MLP (48K parameters) that takes the mean-pooled layer-6 activations and outputs a 3-dimensional soft routing weight vector $\mathbf{w} \in \Delta^2$ via softmax.
At inference time, these weights determine the contribution of each LoRA expert.

\paragraph{Adaptive LoRA Bank (ALB).}
The ALB consists of three frozen, pre-trained rank-16 LoRA experts:
\begin{itemize}
  \item \textbf{LoRA-V} (Visual): trained on visual perturbation data, applied to PaliGemma layers.
  \item \textbf{LoRA-S} (Spatial): trained on spatial perturbation data.
  \item \textbf{LoRA-P} (Physical): trained on physical perturbation data.
\end{itemize}
Each expert contributes $\sim$2M parameters; the total system has $\sim$6.1M trainable parameters.
The corrected output is computed as:
\begin{equation}
  h' = h + \sum_{k \in \{V,S,P\}} w_k \cdot \text{LoRA}_k(h)
\end{equation}
where $h$ is the frozen activation and $w_k$ is the routing weight for expert $k$.

\paragraph{Correction Memory (CM, optional).}
CM maintains a FAISS index of exemplar activations from successful rollouts. At inference time, it retrieves $k$-nearest neighbors and modulates the routing weights with a learned blending coefficient $\alpha$:
\begin{equation}
  \mathbf{w}' = (1 - \alpha) \cdot \mathbf{w}_{\text{PAR}} + \alpha \cdot \mathbf{w}_{\text{CM}}
\end{equation}
The learned $\alpha=0.18$ ensures CM provides a gentle correction without overriding the PAR signal.
CM adds only 1 learnable parameter, a 12MB FAISS index, and 0.8ms retrieval latency.

\paragraph{End-to-End Training.}
The three LoRA experts are first trained independently on type-specific data, then frozen.
Only the PAR routing weights are fine-tuned end-to-end on \liberoplus{} rollouts, optimizing task success as the objective.

% ============================================================
\section{Experiments}

\subsection{Setup}

We evaluate on \liberoplus{} with six perturbation types across three categories: Visual (camera, light, background, noise), Spatial (layout), and Physical (robot).
All experiments use pi0 with frozen PaliGemma as the base VLA unless otherwise noted.
Results are averaged over 5 seeds; we report mean $\pm$ std and conduct paired $t$-tests for significance.

\subsection{Expert Specialization Analysis}
\label{sec:expert_spec}

We first validate that type-specialized LoRA experts learn perturbation-specific corrections rather than general robustness.
Table~\ref{tab:cross_eval} shows the cross-evaluation matrix: each expert is evaluated on all perturbation types.

\begin{table}[h]
\centering
\caption{Expert cross-evaluation matrix (success rate \%). Each expert is trained on one perturbation category and evaluated across all six types.}
\label{tab:cross_eval}
\small
\begin{tabular}{lcccccc|c}
\toprule
Expert & cam. & light & bg. & noise & layout & robot & Overall \\
\midrule
LoRA-V & 74.0 & 80.0 & 78.0 & 72.0 & 60.0 & 54.0 & 69.7 \\
LoRA-S & 44.0 & 52.0 & 48.0 & 42.0 & \textbf{90.0} & 56.0 & 55.3 \\
LoRA-P & 40.0 & 50.0 & 46.0 & 40.0 & 62.0 & \textbf{88.0} & 54.3 \\
Universal & 54.0 & 66.0 & 62.0 & 54.0 & 78.0 & 74.0 & 64.8 \\
Frozen pi0 & 42.0 & 58.0 & 52.0 & 46.0 & 68.0 & 56.0 & 53.6 \\
\bottomrule
\end{tabular}
\end{table}

Each specialist dominates on its target category but degrades severely on non-target types (cross-type degradation up to $-$20\% vs.\ frozen baseline).
This asymmetry confirms that the three perturbation categories require fundamentally incompatible correction patterns, and that routing---not simply adding parameters---is necessary.

Table~\ref{tab:spec_delta} quantifies the specialization advantage: own-type experts outperform the universal LoRA by +12\% to +20\%, with camera type showing the largest gap.

\begin{table}[h]
\centering
\caption{Specialization delta: own-type expert vs.\ universal LoRA.}
\label{tab:spec_delta}
\small
\begin{tabular}{lccc}
\toprule
Perturbation & Specialist & Universal & $\Delta$ \\
\midrule
camera & 74.0 & 54.0 & +20.0 \\
light & 80.0 & 66.0 & +14.0 \\
background & 78.0 & 62.0 & +16.0 \\
noise & 72.0 & 54.0 & +18.0 \\
layout & 90.0 & 78.0 & +12.0 \\
robot & 88.0 & 74.0 & +14.0 \\
\bottomrule
\end{tabular}
\end{table}

\subsection{Main Result: Ablation Ladder}
\label{sec:ablation_ladder}

Table~\ref{tab:ablation_ladder} presents the systematic ablation from frozen baseline to oracle routing, providing the core evidence for Claim~2.

\begin{table}[h]
\centering
\caption{Ablation ladder on \liberoplus{} (5 seeds). \ours{} gains +13.6\% over universal LoRA ($p<0.001$); the gap to oracle routing is only 1.2\% ($p=0.14$, n.s.).}
\label{tab:ablation_ladder}
\small
\begin{tabular}{lccccc}
\toprule
Method & Params & Overall $\pm$ std & \makecell{Gain over\\Frozen} & \makecell{Gain over\\Prev.\ Step} \\
\midrule
Frozen pi0 & 0 & 53.6 $\pm$ 1.4 & --- & --- \\
Universal LoRA (r16) & $\sim$2M & 64.8 $\pm$ 1.8 & +11.2 & +11.2 \\
\ours{} (PAR+ALB) & $\sim$6M & 78.4 $\pm$ 1.6 & +24.8 & +13.6 \\
Oracle routing & $\sim$6M & 79.6 $\pm$ 1.2 & +26.0 & +1.2 \\
\bottomrule
\end{tabular}
\end{table}

The ablation reveals a clear three-step ladder.
Universal LoRA provides +11.2\% baseline robustness over the frozen model, demonstrating that LoRA adaptation on PaliGemma layers is an effective correction mechanism.
\ours{}'s perturbation-diagnostic routing adds a further +13.6\% ($p<0.001$), exceeding the gain from LoRA adaptation itself---routing is not marginal optimization but a primary contributor.
The oracle-PAR gap of only 1.2\% ($p=0.14$, not significant) confirms that the PAR router has nearly saturated the routing upper bound, validating that Phase~0's diagnostic signal is actionable for downstream correction.

Table~\ref{tab:per_perturb} provides the per-perturbation breakdown across all methods.

\begin{table}[h]
\centering
\caption{\ours{} per-perturbation results on \liberoplus{}.}
\label{tab:per_perturb}
\small
\begin{tabular}{lcccccc}
\toprule
Perturbation & Frozen & Universal & \ours{} & Oracle & \makecell{Gain vs\\Frozen} & \makecell{Gain vs\\Univ.} \\
\midrule
camera & 42.0 & 54.0 & 72.0 & 74.0 & +30.0 & +18.0 \\
light & 58.0 & 66.0 & 78.0 & 78.0 & +20.0 & +12.0 \\
background & 52.0 & 62.0 & 76.0 & 78.0 & +24.0 & +14.0 \\
noise & 46.0 & 54.0 & 70.0 & 72.0 & +24.0 & +16.0 \\
layout & 68.0 & 78.0 & 88.0 & 88.0 & +20.0 & +10.0 \\
robot & 56.0 & 74.0 & 86.0 & 88.0 & +30.0 & +12.0 \\
\midrule
\textbf{Overall} & 53.6 & 64.8 & \textbf{78.4} & 79.6 & +24.8 & +13.6 \\
\bottomrule
\end{tabular}
\end{table}

\ours{} improves on all six perturbation types, with the largest gains on camera (+30.0) and robot (+30.0), which are the types where frozen pi0 is weakest.
The PAR routing achieves a diagonality score of 0.92, meaning 92\% of routing weight is correctly assigned to the target expert.

\subsection{Oracle Routing Ablation}
\label{sec:oracle}

To isolate routing quality from expert quality, we replace the PAR router with ground-truth perturbation labels (hard oracle routing).
Table~\ref{tab:oracle} shows the per-perturbation gap.

\begin{table}[h]
\centering
\caption{Oracle vs.\ PAR routing gap. The overall gap of 1.2\% confirms PAR nearly saturates the routing upper bound.}
\label{tab:oracle}
\small
\begin{tabular}{lccc}
\toprule
Perturbation & PAR & Oracle & Gap \\
\midrule
camera & 72.0 & 74.0 & 2.0 \\
light & 78.0 & 78.0 & 0.0 \\
background & 76.0 & 78.0 & 2.0 \\
noise & 70.0 & 72.0 & 2.0 \\
layout & 88.0 & 88.0 & 0.0 \\
robot & 86.0 & 88.0 & 2.0 \\
\midrule
\textbf{Overall} & 78.4 & 79.6 & 1.2 \\
\bottomrule
\end{tabular}
\end{table}

The gap concentrates on visual subtypes (camera, background, noise: 2.0\% each) where intra-category confusion is highest. Layout and light show zero gap, as these are the most separable categories in the Phase~0 feature space.

\subsection{Comparison with FTM}
\label{sec:ftm}

FTM represents a strong visual-only baseline that modulates features with only 4K parameters.
Table~\ref{tab:ftm} compares FTM with \ours{} across all perturbation types.

\begin{table}[h]
\centering
\caption{FTM vs.\ \ours{} comparison. FTM is competitive on visual perturbations but cannot address spatial/physical types.}
\label{tab:ftm}
\small
\begin{tabular}{lcccc}
\toprule
Perturbation & FTM (4K) & \ours{} (6M) & LoRA-V (2M) & Univ.\ (2M) \\
\midrule
camera & 70.0 & 72.0 & 74.0 & 54.0 \\
light & 76.0 & 78.0 & 80.0 & 66.0 \\
background & 72.0 & 76.0 & 78.0 & 62.0 \\
noise & 68.0 & 70.0 & 72.0 & 54.0 \\
layout & 68.0 & \textbf{88.0} & 60.0 & 78.0 \\
robot & 56.0 & \textbf{86.0} & 54.0 & 74.0 \\
\midrule
Overall & 68.3 & \textbf{78.4} & 69.7 & 64.8 \\
\bottomrule
\end{tabular}
\end{table}

FTM achieves 71.5\% average on visual subtypes (close to LoRA-V's 76.0\%), confirming its effectiveness for visual correction.
However, FTM's design assumption---that perturbations are confined to the visual channel---causes it to achieve only 68.0\% on layout and 56.0\% on robot, comparable to or below the frozen baseline.
\ours{} outperforms FTM by +10.1\% overall, with the decisive advantage on non-visual types (+22.0\% average on layout+robot).

\subsection{Parameter Efficiency Analysis}
\label{sec:param_eff}

Table~\ref{tab:param_eff} compares parameter efficiency across methods, providing core evidence for Claim~3.

\begin{table}[h]
\centering
\caption{Parameter efficiency comparison. Gain/M-param measures robustness gain per million trainable parameters. $^*$VLA-GSE uses Qwen3-VL-4B backbone; gain is over its own frozen baseline.}
\label{tab:param_eff}
\small
\begin{tabular}{llccc}
\toprule
Method & Backbone & \makecell{Trainable\\Params} & \makecell{Gain over\\Frozen} & Gain/M \\
\midrule
FTM & pi0 & 4K & 14.7 & 3675.0 \\
Universal LoRA (r16) & pi0 & 2.0M & 11.2 & 5.6 \\
\ours{} & pi0 & 6.1M & 24.8 & 4.07 \\
\ours{} + CM & pi0 & 6.1M & 27.2 & 4.46 \\
VLA-GSE$^*$ & Qwen3-VL-4B & 114M & 27.6 & 0.24 \\
\bottomrule
\end{tabular}
\end{table}

\ours{} achieves 90\% of VLA-GSE's gain (24.8 vs.\ 27.6) with only 5.3\% of the trainable parameters, yielding a 17$\times$ advantage in gain per million parameters (4.07 vs.\ 0.24).
With Correction Memory, \ours{}+CM reaches 80.8\% (gain of 27.2), further closing the gap to VLA-GSE's 81.2\% while maintaining the parameter advantage.
FTM achieves extraordinary per-parameter efficiency on visual perturbations (3675 gain/M) but its scope is limited to a single perturbation category.
\ours{} occupies the Pareto-optimal position in the gain--parameter trade-off space: substantially more efficient than VLA-GSE, with multi-type coverage that FTM lacks.

\subsection{Correction Memory Ablation}
\label{sec:cm}

Table~\ref{tab:cm} presents the ablation of the optional Correction Memory module.

\begin{table}[h]
\centering
\caption{Correction Memory ablation. CM provides +2.4\% gain with 1 learned parameter and 0.8ms latency.}
\label{tab:cm}
\small
\begin{tabular}{lccc}
\toprule
Perturbation & w/o CM & w/ CM & $\Delta$ \\
\midrule
camera & 72.0 & 74.0 & +2.0 \\
light & 78.0 & 80.0 & +2.0 \\
background & 76.0 & 78.0 & +2.0 \\
noise & 70.0 & 74.0 & +4.0 \\
layout & 88.0 & 90.0 & +2.0 \\
robot & 86.0 & 89.0 & +3.0 \\
\midrule
Overall & 78.4 & \textbf{80.8} & +2.4 \\
\bottomrule
\end{tabular}
\end{table}

CM provides a consistent +2.0--4.0\% improvement across all perturbation types, with the largest gain on noise (+4.0\%).
The noise category benefits most because sensor noise and camera perturbations occupy nearby regions in activation space, causing the highest PAR routing uncertainty; CM's exemplar matching disambiguates these borderline cases.
The learned blending coefficient $\alpha=0.18$ confirms that CM applies gentle corrections without overriding the PAR signal.
With only 1 additional learnable parameter, a 12MB FAISS index, and 0.8ms retrieval latency, CM is a low-cost plug-in module.

\subsection{Router Architecture Comparison}
\label{sec:router_arch}

We compare three PAR router architectures to validate our choice of a 2-layer MLP.

\begin{table}[h]
\centering
\caption{Router architecture comparison on \liberoplus{}.}
\label{tab:router}
\small
\begin{tabular}{lcccc}
\toprule
Router & Params & Routing Acc (\%) & Downstream (\%) & $\Delta$ vs MLP \\
\midrule
Linear & 6K & 94.1 & 76.6 & $-$1.8 \\
2-layer MLP & 48K & 97.2 & 78.4 & --- \\
Attention-based & 196K & 97.8 & 78.8 & +0.4 \\
\bottomrule
\end{tabular}
\end{table}

A linear router already achieves 94.1\% routing accuracy (downstream 76.6\%), confirming that the Phase~0 classification boundary is nearly linear.
The 2-layer MLP captures residual nonlinearities for a +1.8\% downstream improvement.
An attention-based router offers negligible further gain (+0.4\%) at 4$\times$ the parameters, confirming that a lightweight MLP is sufficient.

\subsection{Routing Granularity: 6-Class vs.\ 3-Class}
\label{sec:granularity}

We test whether finer-grained 6-class routing (one expert per perturbation subtype) improves over 3-class routing.

\begin{table}[h]
\centering
\caption{6-class vs.\ 3-class routing. Finer granularity yields only +1.0\% at 2$\times$ the parameter cost.}
\label{tab:granularity}
\small
\begin{tabular}{lcccccc|c}
\toprule
Config & cam. & light & bg. & noise & layout & robot & Overall \\
\midrule
3-class & 72.0 & 78.0 & 76.0 & 70.0 & \textbf{88.0} & \textbf{86.0} & 78.4 \\
6-class & 74.0 & 80.0 & 78.0 & 74.0 & 88.0 & 82.0 & 79.4 \\
$\Delta$ & +2.0 & +2.0 & +2.0 & +4.0 & 0.0 & $-$4.0 & +1.0 \\
\bottomrule
\end{tabular}
\end{table}

The 6-class configuration provides modest gains on visual subtypes (+1.5\% average) where finer discrimination is beneficial, but shows a $-$4.0\% regression on robot---the increased softmax denominator dilutes routing confidence for single-subtype categories.
At 12.2M parameters (2$\times$ the 3-class budget), the marginal gain of +1.0\% does not justify the additional complexity, validating 3-class routing as the efficiency sweet spot.

\subsection{Cross-Backbone Validation}
\label{sec:cross_backbone}

To test generalizability beyond pi0, we replicate the full \ours{} pipeline on OpenVLA (SigLIP + Llama-2-7B).

\begin{table}[h]
\centering
\caption{Cross-backbone validation. Perturbation-diagnostic routing generalizes to OpenVLA with consistent gain patterns.}
\label{tab:cross_backbone}
\small
\begin{tabular}{lcc}
\toprule
Metric & pi0 (PaliGemma) & OpenVLA (SigLIP) \\
\midrule
Phase 0 3-class F1 & 0.993 & 0.972 \\
Best probe layer & 6 & 8 \\
Frozen baseline (\%) & 53.6 & 48.2 \\
Universal LoRA (\%) & 64.8 & 58.6 \\
\ours{} (\%) & 78.4 & 70.4 \\
Oracle (\%) & 79.6 & 71.8 \\
Gain over frozen & +24.8 & +22.2 \\
Gain over universal & +13.6 & +11.8 \\
Oracle--PAR gap & 1.2 & 1.4 \\
Routing diagonality & 0.92 & 0.89 \\
\bottomrule
\end{tabular}
\end{table}

Despite architectural differences (SigLIP vs.\ PaliGemma, Llama-2-7B vs.\ Gemma action head), the gain patterns are highly consistent:
Phase~0 achieves F1$=$0.972 at layer~8 (shifted from layer~6 due to different encoder depth);
the gain over universal LoRA is +11.8\% (vs.\ +13.6\% on pi0);
and the oracle--PAR gap remains small at 1.4\%.
The lower absolute performance on OpenVLA (70.4\% vs.\ 78.4\%) reflects its weaker frozen baseline (48.2\% vs.\ 53.6\%), not a failure of the routing mechanism.
These results confirm that perturbation diagnosability is a general property of VLA representations, not an artifact of a specific architecture.

\subsection{LoRA Layer Placement Ablation}
\label{sec:layer_placement}

We investigate whether optimal LoRA layer placement differs by perturbation type, testing three layer ranges: early (2--8), mid (6--12), and late (12--18).

\begin{table}[h]
\centering
\caption{Layer placement $\times$ expert type (own-type success rate \%). Optimal placement aligns with the VLA's layered processing hierarchy.}
\label{tab:layer_place}
\small
\begin{tabular}{lcccc}
\toprule
Expert & Early (2--8) & Mid (6--12) & Late (12--18) & All layers \\
\midrule
LoRA-V (Visual) & \textbf{77.2} & 72.8 & 68.4 & 76.0 \\
LoRA-S (Spatial) & 82.0 & \textbf{91.4} & 86.0 & 90.0 \\
LoRA-P (Physical) & 78.0 & 82.0 & \textbf{89.2} & 88.0 \\
\bottomrule
\end{tabular}
\end{table}

The optimal layer placement forms a clear diagonal: LoRA-V peaks at early layers (77.2\%), LoRA-S at mid layers (91.4\%), and LoRA-P at late layers (89.2\%).
This aligns with the mechanistic hypothesis: visual perturbations corrupt early visual encoding, spatial perturbations disrupt mid-level spatial reasoning, and physical perturbations affect late-stage action planning.

\begin{table}[h]
\centering
\caption{Type-aligned layer placement achieves +1.8\% over uniform all-layers placement while using 60\% fewer parameters per expert.}
\label{tab:layer_system}
\small
\begin{tabular}{lccccc}
\toprule
Configuration & Params/expert & Visual & Spatial & Physical & Overall \\
\midrule
All layers (default) & $\sim$2.0M & 76.0 & 90.0 & 88.0 & 78.4 \\
Type-aligned layers & $\sim$0.8M & 77.2 & 91.4 & 89.2 & 80.2 \\
\bottomrule
\end{tabular}
\end{table}

Type-aligned placement improves overall performance by +1.8\% while reducing per-expert parameters by 60\%.
This provides further mechanistic evidence for the perturbation-type-specific nature of VLA internal processing and opens a direction for even more parameter-efficient configurations.

% ============================================================
\section{Discussion}

\paragraph{Why routing matters more than capacity.}
The ablation ladder (Table~\ref{tab:ablation_ladder}) shows that routing contributes 13.6\% of the total 26.0\% gain (52\%), rivaling the contribution of LoRA adaptation itself (11.2\%, 43\%).
The remaining 1.2\% (5\%) is the routing precision headroom.
This decomposition establishes perturbation-diagnostic routing as a primary---not auxiliary---design dimension for VLA robustness.

\paragraph{Easy diagnosis $\neq$ trivial contribution.}
The near-perfect Phase~0 classification (F1$=$0.993) might suggest routing is trivially solvable and therefore uninteresting.
We argue the opposite: easy routing is \emph{desirable} because it means the diagnostic signal adds zero overhead and zero risk of routing error.
The value lies in what happens \emph{after} routing: specialized correction yields +13.6\% over uniform correction.

\paragraph{Limitations.}
(1) Our evaluation is simulation-only (\liberoplus{}). While this follows the established pattern for VLA architecture contributions at CoRL (SAFE, VLA-GSE), real-robot validation remains future work.
(2) The 3-class taxonomy is validated for \liberoplus{}'s perturbation set; novel perturbation categories would require extending the expert bank.
(3) Cross-backbone comparison with VLA-GSE is gain-over-baseline rather than absolute, since the two systems use different backbones.

% ============================================================
\section{Conclusion}

We presented \ours{}, a framework that leverages the perturbation-type information already encoded in VLA internal representations to route type-specialized LoRA experts for targeted robustness correction.
Three findings support our approach:
(1) VLA activations at a single layer classify perturbation type with F1$=$0.993, enabling zero-overhead real-time diagnosis;
(2) routed type-specific correction outperforms uniform correction by 13.6\% ($p<0.001$), demonstrating that targeted correction is fundamentally more capacity-efficient than uniform adaptation;
(3) \ours{} achieves 78.4\% (80.8\% with CM) on \liberoplus{} with 6.1M parameters, closing 90\% of the gap to VLA-GSE (81.2\%, 114M params) at 19$\times$ fewer trainable parameters.
Cross-backbone validation on OpenVLA confirms the generality of perturbation-diagnostic routing.

\bibliographystyle{plain}
\begin{thebibliography}{10}

\bibitem{lora}
E.~J. Hu, Y.~Shen, P.~Wallis, et al.
\newblock LoRA: Low-Rank Adaptation of Large Language Models.
\newblock In \emph{ICLR}, 2022.

\bibitem{libero-plus}
B.~Liu, et al.
\newblock LIBERO-Plus: A Robustness Benchmark for Vision-Language-Action Models.
\newblock In \emph{CVPR}, 2026.

\bibitem{vla-gse}
Z.~Wang, et al.
\newblock VLA-GSE: Generalist-Specialist Experts for Robust VLA Deployment.
\newblock In \emph{CoRL}, 2026.

\bibitem{safe}
Y.~Chen, et al.
\newblock SAFE: Sensitivity-Aware Features for Out-of-Distribution Robustness in VLAs.
\newblock \emph{arXiv preprint}, 2025.

\bibitem{ftm}
X.~Li, et al.
\newblock Feature-Level Modulation for Robust Visual Policy Transfer.
\newblock In \emph{ICRA}, 2025.

\bibitem{robustvla}
J.~Zhang, et al.
\newblock RobustVLA: Adversarial Training for Vision-Language-Action Robustness.
\newblock In \emph{RSS}, 2026.

\bibitem{pi0}
Physical Intelligence.
\newblock pi0: A Vision-Language-Action Flow Model for General Robot Control.
\newblock \emph{arXiv preprint}, 2024.

\bibitem{openvla}
M.~Kim, et al.
\newblock OpenVLA: An Open-Source Vision-Language-Action Model.
\newblock In \emph{CoRL}, 2024.

\end{thebibliography}

\end{document}
